% problem_id: S001-spaces-perfect-proposition-detecting-bounded-above
% source_id: S001 (Stacks Project)
% source_file: spaces-perfect.tex
% statement_locator: spaces-perfect.tex lines 7081--7105
% proof_locator: spaces-perfect.tex lines 7107--7202
% env: proposition
% pairing: tight (verified)
% status: complete (statement + author proof verbatim + reason + locators)
%
\begin{proposition}
\label{proposition-detecting-bounded-above}
Let $S$ be a scheme.
Let $X$ be a quasi-compact and quasi-separated algebraic space over $S$.
Let $G \in D_{perf}(\mathcal{O}_X)$ be a perfect complex which generates
$D_\QCoh (\mathcal{O}_X)$. Let $E \in D_\QCoh (\mathcal{O}_X)$.
The following are equivalent
\begin{enumerate}
\item $E \in D^-_\QCoh (\mathcal{O}_X)$,
\item $\Hom_{D(\mathcal{O}_X)}(G[-i], E) = 0$ for $i \gg 0$,
\item $\Ext^i_X(G, E) = 0$ for $i \gg 0$,
\item $R\Hom_X(G, E)$ is in $D^-(\mathbf{Z})$,
\item $H^i(X, G^\vee \otimes_{\mathcal{O}_X}^\mathbf{L} E) = 0$
for $i \gg 0$,
\item $R\Gamma(X, G^\vee \otimes_{\mathcal{O}_X}^\mathbf{L} E)$
is in $D^-(\mathbf{Z})$,
\item for every perfect object $P$ of $D(\mathcal{O}_X)$
\begin{enumerate}
\item the assertions (2), (3), (4) hold with $G$ replaced by $P$, and
\item $H^i(X, P \otimes_{\mathcal{O}_X}^\mathbf{L} E) = 0$ for $i \gg 0$,
\item $R\Gamma(X, P \otimes_{\mathcal{O}_X}^\mathbf{L} E)$
is in $D^-(\mathbf{Z})$.
\end{enumerate}
\end{enumerate}
\end{proposition}

\begin{proof}
Assume (1). Since
$\Hom_{D(\mathcal{O}_X)}(G[-i], E) = \Hom_{D(\mathcal{O}_X)}(G, E[i])$
we see that this is zero for $i \gg 0$ by
Lemma \ref{lemma-ext-from-perfect-into-bounded-QCoh}. This proves
that (1) implies (2).

\medskip\noindent
Parts (2), (3), (4) are equivalent by the discussion in
Cohomology on Sites, Section \ref{sites-cohomology-section-global-RHom}.
Part (5) and (6) are equivalent as $H^i(X, -) = H^i(R\Gamma(X, -))$
by definition. The equivalent conditions (2), (3), (4) are
equivalent to the equivalent conditions (5), (6) by
Cohomology on Sites, Lemma \ref{sites-cohomology-lemma-dual-perfect-complex}
and the fact that $(G[-i])^\vee = G^\vee[i]$.

\medskip\noindent
It is clear that (7) implies (2). Conversely,
let us prove that the equivalent conditions (2) -- (6) imply (7).
Recall that $G$ is a  classical generator for $D_{perf}(\mathcal{O}_X)$ by
Remark \ref{remark-classical-generator}.
For $P \in D_{perf}(\mathcal{O}_X)$ let $T(P)$ be the assertion that
$R\Hom_X(P, E)$ is in $D^-(\mathbf{Z})$.
Clearly, $T$ is inherited by direct sums,
satisfies the 2-out-of-three property for distinguished
triangles, is inherited by direct summands, and is preserved by shifts.
Hence by Derived Categories, Remark \ref{derived-remark-check-on-generator}
we see that (4) implies $T$ holds on all of $D_{perf}(\mathcal{O}_X)$.
The same argument works for all other properties, except that for property
(7)(b) and (7)(c) we also use that $P \mapsto P^\vee$ is a self
equivalence of $D_{perf}(\mathcal{O}_X)$. Small detail omitted.

\medskip\noindent
We will prove the equivalent conditions (2) -- (7) imply (1)
using the induction principle of
Lemma \ref{lemma-induction-principle}.

\medskip\noindent
First, we prove (2) -- (7) $\Rightarrow$ (1) if $X$ is affine.
This follows from the case of schemes, see
Derived Categories of Schemes, Proposition
\ref{perfect-proposition-detecting-bounded-above}.

\medskip\noindent
Now assume $(U \subset X, j : V \to X)$ is an elementary distinguished
square of quasi-compact and quasi-separated algebraic spaces over $S$
and assume the implication (2) -- (7) $\Rightarrow$ (1)
is known for $U$, $V$, and $U \times_X V$. To finish the proof
we have to show the implication (2) -- (7) $\Rightarrow$ (1) for $X$.
Suppose $E \in D_\QCoh(\mathcal{O}_X)$ satisfies (2) -- (7).
By Lemma \ref{lemma-direct-summand-of-a-restriction} and
Theorem \ref{theorem-bondal-van-den-Bergh} there exists a perfect
complex $Q$ on $X$ such that $Q|_U$ generates $D_\QCoh (\mathcal{O}_U)$.

\medskip\noindent
Say $V = \Spec(A)$. Let $Z \subset V$ be the reduced closed subscheme
which is the inverse image of $X \setminus U$ and maps isomorphically to it
(see Definition \ref{definition-elementary-distinguished-square}).
This is a retrocompact closed subset of $V$.
Choose $f_1, \ldots, f_r \in A$ such that
$Z = V(f_1, \ldots, f_r)$. Let $K \in D(\mathcal{O}_V)$ be the perfect
object corresponding to the Koszul complex on $f_1, \ldots, f_r$ over $A$.
Note that since $K$ is supported on $Z$, the pushforward
$K' = Rj_*K$ is a perfect object of $D(\mathcal{O}_X)$ whose
restriction to $V$ is $K$ (see Lemmas \ref{lemma-pushforward-perfect}
and \ref{lemma-pushforward-with-support-in-open}).
By assumption, we know $R\Hom_{\mathcal{O}_X}(Q, E)$ and
$R\Hom_{\mathcal{O}_X}(K', E)$ are bounded above.

\medskip\noindent
By Lemma \ref{lemma-pushforward-with-support-in-open}
we have $K' =  j_!K$ and hence
$$
\Hom_{D(\mathcal{O}_X)}(K'[-i], E) = \Hom_{D(\mathcal{O}_V)}(K[-i], E|_V) = 0
$$
for $i \gg 0$. Therefore, we may apply
Derived Categories of Schemes,
Lemma \ref{perfect-lemma-orthogonal-koszul-first-variant} to $E|_V$ to
obtain an integer $a$ such that
$\tau_{\geq a}(E|_V) =
\tau_{\geq a} R (U \times_X V \to V)_* (E|_{U \times_X V})$.
Then $\tau_{\geq a} E = \tau_{\geq a} R (U \to X)_* (E |_U)$
(check that the canonical map is an isomorphism after restricting to
$U$ and to $V$). Hence using Lemma \ref{lemma-bounded-truncation}
twice we see that
$$
\Hom_{D(\mathcal{O}_U)}(Q|_U [-i], E|_U) =
\Hom_{D(\mathcal{O}_X)}(Q[-i], R (U \to X)_* (E|_U)) = 0
$$
for $i \gg 0$. Since the Proposition holds for $U$ and the generator
$Q|_U$, we have $E|_U \in D^-_\QCoh(\mathcal{O}_U)$. But then since
the functor $R (U \to X)_*$ preserves $D^-_\QCoh$
(by Lemma \ref{lemma-quasi-coherence-direct-image}), we get
$\tau_{\geq a}E \in D^-_\QCoh(\mathcal{O}_X)$. Thus
$E \in D^-_\QCoh (\mathcal{O}_X)$.
\end{proof}
